% =============================================================================
%  90_appendix.tex —— 附录 A 术语表、B 校勘记、C 索引、D 精读指南
% =============================================================================
\section{术语对照表}\label{app:terms}

下表列出正文出现的术语，按主题分组；「首次」一栏给出首次出现的位置。
同一英文词在正文中有不同译法时，以本表为准。

{\small
\renewcommand{\arraystretch}{1.28}
\begin{longtable}{@{}p{4.2cm}p{5.0cm}p{4.8cm}@{}}
\toprule
中文 & 英文 & 首次出现／备注\\
\midrule
\endfirsthead
\toprule
中文 & 英文 & 首次出现／备注\\
\midrule
\endhead
\bottomrule
\endlastfoot
\multicolumn{3}{@{}l}{\textbf{一、范畴与同调代数}}\\
加性范畴 & additive category & \S\ref{sec:cat}，定义 \ref{def:additive}\\
阿贝尔范畴 & abelian category & 同上\\
对象／态射 & object / morphism & 同上\\
核／余核 & kernel / cokernel & 同上\\
单态射／满态射 & monomorphism / epimorphism & 同上\\
零对象 & zero object & 同上\\
直和 & direct sum & 同上\\
正合（左、右正合） & exact (left, right exact) & \S\ref{sec:cat}\\
链复形 & chain complex & \S\ref{sec:cx}\\
链映射 & chain map & \S\ref{sec:cx}\\
上同调／同调 & cohomology / homology & \S\ref{sec:cx}\\
边界算子 & boundary operator & \S\ref{sec:cx}\\
拟同构 & quasi-isomorphism & \S\ref{sec:qi}\\
局部化 & localisation & \S\ref{sec:cat}\\
屋顶 & roof & \S\ref{sec:cat}\\
分数演算 & calculus of fractions & 概念卡片，\S\ref{sec:cat}\\
导出范畴 & derived category & \S\ref{sec:cat}，定义 \ref{def:db}\\
有界导出范畴 & bounded derived category & \S\ref{sec:cat}，定义 \ref{def:db}\\
有界（下方、上方） & bounded (below, above) & 注 \ref{rem:four}\\
平移函子 & shift functor & \S\ref{sec:cone}\\
特异三角 & distinguished triangle & \S\ref{sec:cone}\\
三角范畴 & triangulated category & \S\ref{sec:cone}\\
八面体引理 & octahedral lemma & \S\ref{sec:cone}\\
五引理 & five-lemma & \S\ref{sec:cone}\\
锥 & cone & \S\ref{sec:cone}，图 \ref{fig:cone}\\
柱 & cylinder & \S\ref{sec:cone}\\
悬挂 & suspension & \S\ref{sec:cone}\\
环路空间 & loop space & \S\ref{sec:cone}\\
谱（稳定同伦） & spectra & \S\ref{sec:cone}\\
同伦 & homotopy & \S\ref{sec:cx}\\
\midrule
\multicolumn{3}{@{}l}{\textbf{二、导出函子}}\\
导出函子 & derived functor & \S\ref{sec:df}\\
右导出函子 & right derived functor & \S\ref{sec:df}\\
左导出函子 & left derived functor & \S\ref{sec:df}\\
适应（类） & adapted (class) & \S\ref{sec:df}，定义 \ref{def:adapted}\\
无环复形 & acyclic complex & 同上\\
内射／投射／平坦 & injective / projective / flat & 同上\\
消解 & resolution & \S\ref{sec:cx}\\
层上同调 & sheaf cohomology & \S\ref{sec:df}\\
超上同调 & hypercohomology & \S\ref{sec:df}\\
谱序列 & spectral sequence & \S\ref{sec:df}\\
细层 & fine sheaf & 例 \ref{eg:derham}\\
单位分解 & partition of unity & 同上\\
完美复形 & perfect complex & 校勘 \ref{kanwu:perfect}\\
对偶复形 & dual complex & 例 \ref{eg:duality}\\
\midrule
\multicolumn{3}{@{}l}{\textbf{三、拓扑}}\\
单纯复形 & simplicial complex & \S\ref{sec:cx}\\
单纯映射 & simplicial map & \S\ref{sec:cx}\\
三角剖分 & triangulation & \S\ref{sec:cx}\\
细分 & subdivision & \S\ref{sec:cx}\\
单纯逼近 & simplicial approximation & \S\ref{sec:cx}\\
单形 & simplex & \S\ref{sec:cx}\\
纤维化／上纤维化 & fibration / cofibration & \S\ref{sec:cone}\\
弱同伦等价 & weak homotopy equivalence & \S\ref{sec:cone}\\
\midrule
\multicolumn{3}{@{}l}{\textbf{四、代数几何}}\\
层 & sheaf & \S\ref{sec:cx}\\
结构层 & structure sheaf & \S\ref{sec:cx}\\
凝聚层 & coherent sheaf & \S\ref{sec:cx}\\
拟凝聚层 & quasi-coherent sheaf & \S\ref{sec:df}\\
局部自由 & locally free & \S\ref{sec:cx}\\
向量丛 & vector bundle & \S\ref{sec:cx}\\
线丛 & line bundle & \S\ref{sec:cx}\\
截面 & section & \S\ref{sec:cx}\\
除子 & divisor & \S\ref{sec:cx}\\
横截 & transverse & \S\ref{sec:cx}\\
正则截面 & regular section & \S\ref{sec:cx}\\
余维 & codimension & \S\ref{sec:cx}\\
内乘 & interior product & 译注 \ref{yizhu:koszul}\\
Koszul 复形 & Koszul complex & \S\ref{sec:cx}\\
挠层 & torsion sheaf & \S\ref{sec:cx}\\
自相交 & self-intersection & \S\ref{sec:cx}\\
陈类 & Chern class & \S\ref{sec:cx}\\
概形 & scheme & \S\ref{sec:cx}\\
摩天层 & skyscraper sheaf & 校勘 \ref{kanwu:perfect}\\
外微分 & exterior derivative & 例 \ref{eg:derham}\\
全纯 de Rham 复形 & holomorphic de Rham complex & 校勘 \ref{kanwu:dolbeault}\\
Dolbeault 复形 & Dolbeault complex & 同上\\
Hodge 群 & Hodge group & 例 \ref{eg:derham}\\
Kähler 流形 & Kähler manifold & 同上\\
\end{longtable}
}

\section{校勘记}\label{app:errata}

这一节逐条记录原文（arXiv:math/0001045）中判定为疑误或表述不周的地方。
判定原则是：先核对数学内容，再核对记号与引文；判不了的标“存疑”，不擅自改动。
正文一律保留原文的数学陈述，校勘以红框 \kanwuRef{n} 就地标出，详情在此。
每条给出位置、原文、判定、依据与建议；建议只供读者参考，不代替原文。

\subsection*{一、确认为误}

\paragraph{校勘 1（节号错位）}
\textbf{位置}：§\ref{sec:intro} 第一段。\textbf{原文}：“Section 1 begins with
chain complexes \ldots{} the subject of Section 2 \ldots{} in Section 3,
leading to cones \ldots{} in Section 4. Finally Section 5 shows \ldots{}”
\textbf{判定}：错。引言之后正文的实际节号是 §2 链复形、§3 拟同构、§4 范畴论、
§5 锥与三角、§6 导出函子。作者写这段话时按“从链复形数起”编号，未把引言计入。
\textbf{依据}：原文自身的 \verb|\label{cx}|、\verb|\label{qi}|、\verb|\label{df}|
交叉引用（\verb|\ref| 编译结果为 2、3、6），以及全文目录。
\textbf{建议}：正文改用实际节号（本书已改），或在引言中把“Section $n$”整体减一。

\paragraph{校勘 2（“末三个箭头”应为“末三项”）}
\textbf{位置}：§\ref{sec:cx} 末段。\textbf{原文}：“tensoring the exact sequence
(\ref{eq:div}) with $\OO_{D_2}$ gives a sequence in which the final three
arrows are still exact.”\textbf{判定}：笔误。正合列 \eqref{eq:div} 共有三个箭头；
张量积后仍正合的是末\emph{三}项（两条箭头）。\textbf{依据}：右正合性的定义。
\textbf{建议}：改为“末三项”。

\paragraph{校勘 3（§3“同调相同”的例子不成立）}
\textbf{位置}：§\ref{sec:qi} 第二段。\textbf{原文}：复形
$\C[x,y]^{\oplus 2}\xrightarrow{(x,y)}\C[x,y]$ 与
$\C[x,y]\xrightarrow{0}\C$“have the same homology but are not
quasi-isomorphic”。\textbf{判定}：前半句错。前者是 $\C$ 的 Koszul 消解，
$H_0=\C$、$H_1=0$；后者 $H_1=\ker 0=\C[x,y]\ne0$、$H_0=\C$。两者不同调。
\textbf{依据}：直接计算；后者在 $1$ 次的上同调是整个多项式环。
\textbf{建议}：改用标准反例：$A=\C[\epsilon]/(\epsilon^2)$，
$K=\{A\xrightarrow{\epsilon}A\}$（$H_0=H_1=\C$）与 $\C\oplus\C[1]$
（零微分，$H_0=H_1=\C$）同调相同但不拟同构——因为
$\Hom_A(\C,A)=0$，$H_1$ 上的同构无法由链映射实现。详见译注
\ref{yizhu:homology}。

\paragraph{校勘 4（“Dolbeault 复形”名实不符；Hodge 指标互换）}
\textbf{位置}：§\ref{sec:df} 例 \ref{eg:derham}。\textbf{原文}：把
$\C\to\{\OO\xrightarrow{\partial}\Omega^{1,0}\xrightarrow{\partial}\ldots
\to\Omega^{n,0}\}$ 称为“the Dolbeault complex or resolution”，并写
$\bigoplus H^i(\Omega^{j,0})=\bigoplus H^{i,j}(X)$。
\textbf{判定}：名实不符。带全纯外微分 $\partial$ 的是\emph{全纯 de Rham 复形}；
Dolbeault 复形是 $(\Omega^{0,\bullet},\bar\partial)$。另外
$H^{i,j}=H^i(\Omega^{j,0})$ 与通行约定 $H^{p,q}=H^q(X,\Omega^p)$ 指标互换。
\textbf{依据}：Griffiths--Harris《Principles of Algebraic Geometry》第 0 章、
第 3 章；Hartshorne《Algebraic Geometry》II.7 与附录。
\textbf{建议}：正文改称“全纯 de Rham 复形”，Hodge 群写成
$H^{i,j}=H^j(\Omega^i)$；若要保留 Dolbeault 一词，则复形应换成
$(\Omega^{0,\bullet},\bar\partial)$。

\paragraph{校勘 5（“内射层 $\subset$ 拟凝聚层”表述不确）}
\textbf{位置}：§\ref{sec:df} 定义 \ref{def:adapted} 之后。\textbf{原文}：
“Examples of such $\RR\subset\A$ include $\{$injective sheaves$\}\subset
\{$quasi-coherent sheaves$\}$”。\textbf{判定}：不确。层范畴的内射层一般并不
拟凝聚（如 $\mathbb Q/\Z$）。作者本意是拟凝聚层范畴里的内射对象。
\textbf{依据}：Hartshorne《Residues and Duality》II.7（Noether 概形上
$\mathrm{QCoh}$ 有足够内射对象）。
\textbf{建议}：改写为“$\{$拟凝聚层中的内射对象$\}\subset$ 拟凝聚层”。

\paragraph{校勘 6（Leray 与局部到整体谱序列的记号残缺）}
\textbf{位置}：§\ref{sec:df} 例 \ref{eg:composition}。\textbf{原文}：
“$H^{i+j}_X\Leftarrow H^i_Y(R^jp_*)$”与
“$\Ext^{i+j}\Leftarrow H^i(\mathcal Ext^j)$”。\textbf{判定}：记号残缺，
缺少 $E_2$ 页的指标与收敛对象。\textbf{依据}：标准形式
$E_2^{i,j}=H^i(Y,R^jp_*)\Rightarrow H^{i+j}(X)$，
$E_2^{i,j}=H^i(X,\mathcal Ext^j)\Rightarrow\Ext^{i+j}(A,B)$
（Weibel 5.6；Hartshorne《Residues and Duality》II.7）。
\textbf{建议}：按上述完整形式书写。

\paragraph{校勘 7（对偶公式缺“完美复形”条件）}
\textbf{位置}：§\ref{sec:df} 例 \ref{eg:duality}。\textbf{原文}：
“$\R\mathcal Hom(A,B)\cong B\Ltimes A\ve$”，其中
$A\ve=\R\mathcal Hom(A,\OO)$。\textbf{判定}：缺条件。该式的标准形式要求
$A$ 是完美复形（此时由求值映射 $A\ve\Ltimes A\to\OO$ 给出同构），
原文对一般的 $A$、$B$ 断言同构，条件不成立。\textbf{依据}：取 $X$ 为光滑曲线、
$x\in X$、$A=B=k_x$（$x$ 处的摩天层）。$\OO_X$ 无挠，故
$\SheafExt^0(k_x,\OO)=0$，而曲线上 $\SheafExt^1(k_x,\OO)=k_x$，于是
$A\ve=\R\mathcal Hom(k_x,\OO)\cong k_x[1]$；又
$k_x\Ltimes k_x\cong k\oplus k[1]$（$\Tor_0=\Tor_1=k$），故右边
$B\Ltimes A\ve\cong(k\oplus k[1])[1]=k[1]\oplus k[2]$。左边
$\R\mathcal Hom(k_x,k_x)\cong k\oplus k[1]$（$\SheafExt^0=\SheafExt^1=k_x$），
两者差一个平移，不同构。
\textbf{建议}：补上“设 $A$ 为完美复形”。

\paragraph{校勘 8（拼写：“rather then”应为“rather than”）}
\textbf{位置}：§\ref{sec:cone} 第二段。\textbf{原文}：“For some \emph{fixed}
maps $f\colon X\to Y$, rather then equivalence classes of homotopic maps
\ldots”。\textbf{判定}：拼写错误，不影响数学。
\textbf{建议}：改为 rather than。

\subsection*{二、存疑（不改，仅标出）}

\paragraph{存疑 1（$\R\Gamma=\R H\udot$ 的记号）}
§\ref{sec:df} 说“$\R\Gamma$ 的上同调就是标准层上同调 $H^i$”，紧接着写
“$\R\Gamma=\R H\udot$”。$H\udot$ 在全文未定义，疑为“层的复形的整体截面”
或“$H^0$”的笔误。按上下文，$\R\Gamma=\Gamma\circ(\text{内射消解})$ 即可，
不影响结论。读者可忽略此等式。

\paragraph{存疑 2（“第 $i$ 个上同调 $H_i$”）}
§\ref{sec:cone} 说“Taking the $i$th cohomology $H_i$ of each term”，
但 $H_i$ 按全文约定是同调。此处应读作“同调”。译文按同调处理，不改原文符号。

\section{定理、定义与例索引}\label{app:index}

{\small
\renewcommand{\arraystretch}{1.25}
\begin{longtable}{@{}p{1.6cm}p{5.4cm}p{6.4cm}@{}}
\toprule
编号 & 内容 & 位置\\
\midrule
\endfirsthead
\toprule
编号 & 内容 & 位置\\
\midrule
\endhead
\bottomrule
\endlastfoot
定义 \ref{def:additive} & 加性范畴、阿贝尔范畴 & \S\ref{sec:cat}\\
定义 \ref{def:db} & 有界导出范畴 $D^b(\A)$ & \S\ref{sec:cat}\\
注 \ref{rem:four} & 四种导出范畴 $D^\pm$、$D^b$、$D$ & \S\ref{sec:cat}\\
定义 \ref{def:adapted} & 适应类；$K^+(\RR)$ 与 $D^+(\A)$ 等价 & \S\ref{sec:df}\\
例 \ref{eg:tor} & $\Tor_i$ 的对称性 & \S\ref{sec:df}\\
例 \ref{eg:rhom} & $\R\mathcal Hom$ 的两种消解 & \S\ref{sec:df}\\
例 \ref{eg:composition} & Grothendieck、Leray、局部到整体、超上同调谱序列 & \S\ref{sec:df}\\
例 \ref{eg:duality} & 对偶公式 $B\Ltimes A\ve$（需完美复形） & \S\ref{sec:df}\\
例 \ref{eg:derham} & de Rham 定理与 Hodge 谱序列 & \S\ref{sec:df}\\
\midrule
式 \eqref{eq:co} & 单纯链复形与上链复形 & \S\ref{sec:cx}\\
式 \eqref{eq:div} & 除子的正合列 & \S\ref{sec:cx}\\
式 \eqref{eq:wh} & Whitehead 定理的图表 & \S\ref{sec:cx}\\
式 \eqref{eq:dot} & 拟同构的“虚线箭头” & \S\ref{sec:qi}\\
式 \eqref{eq:sp} & $X\to Y\to Y/X\to\Sigma X\to\cdots$ & \S\ref{sec:cone}\\
式 \eqref{eq:les} & 对子的同调长正合列 & \S\ref{sec:cone}\\
图 \ref{fig:cone} & 锥 $C_f$ & \S\ref{sec:cone}\\
图 \ref{fig:doldpuppe} & Dold--Puppe 序列 & \S\ref{sec:cone}\\
表 \ref{tab:intersection} & 除子、截面与相交数的对照 & \S\ref{sec:cx}\\
\end{longtable}
}

\section{精读指南：一张地图、五个论证、五道习题}\label{app:reading}

\subsection*{D.1 全文的论证地图}

全文只有一条主线，可以画成这样：

\begin{center}
\small
\begin{tabular}{@{}c@{}}
\textbf{问题}：想用复形代替（上）同调，但复形的选取不典范\\
$\Downarrow$\\
\textbf{第 2 节}：为什么必须用复形（拓扑的 Whitehead、代数几何的消解、相交理论）\\
$\Downarrow$\\
\textbf{第 3 节}：正确的等同关系是拟同构，不是“同调相同”\\
$\Downarrow$\\
\textbf{第 4 节}：把拟同构可逆化——导出范畴 $D(\A)$ 的定义（范畴论的语言）\\
$\Downarrow$\\
\textbf{第 5 节}：$D(\A)$ 不是阿贝尔范畴；核与余核被锥与特异三角取代\\
$\Downarrow$\\
\textbf{第 6 节}：导出函子 = 在适应类上作消解后再作用 $F$；Tor、Ext、超上同调、
de Rham 都是它的特例
\end{tabular}
\end{center}

一句话概括：\emph{复形好，（上）同调坏；把拟同构变成同构，坏处就消失了。}

\subsection*{D.2 逐节导读}

\paragraph{§\ref{sec:intro} 引言。}读法：把这一节当作“预告片”，不必抠细节。
要记住三件事：作者是给物理学家讲的；Kontsevich 的镜像对称把三角范畴带进了主流；
技术细节查 Gelfand--Manin。注意原文的节号错位（校勘 1）。

\paragraph{§\ref{sec:cx} 链复形。}三个例子层层递进：
(1) 拓扑——$H_*$ 丢 Massey 积，$C\ldot$ 由 Whitehead 定理决定同伦型；
(2) 代数几何——消解把挠层换成局部自由层；
(3) 相交理论——$\OO_D\otimes\OO_D$ 只在横截时对，而复形总对，多出的核是
$\Tor_1$。读的时候把 (3) 的表格亲手算一遍：$D_1=D_2=D$ 时，
$\{L^{-1}|_D\xrightarrow{0}\OO_D\}$ 的 $H_0=\OO_D$、$H_1=L^{-1}|_D$，
而 $c_1(L^{-1}|_D)=-D\cdot D$，这就是自相交数的来源。

\paragraph{§\ref{sec:qi} 拟同构。}核心是“屋顶”与“同调相同 $\ne$ 拟同构”。
原文的例子有误（校勘 3），务必用译注 \ref{yizhu:homology} 的
$k[\epsilon]/\epsilon^2$ 例子替换着读，否则会得到错误印象。

\paragraph{§\ref{sec:cat} 范畴论。}这一节最“干”，但只有两个定义：
加性/阿贝尔范畴（定义 \ref{def:additive}）与 $D^b(\A)$
（定义 \ref{def:db}）。后者要读出三层意思：对象是有界复形；态射是链映射与
拟同构的逆的形式词；形式词之间要作商。屋顶图是把这三层意思画在一张图上。

\paragraph{§\ref{sec:cone} 锥与三角。}全篇的几何心脏。建议按图
\ref{fig:cone}、\ref{fig:doldpuppe} 把拓扑的锥、Dold--Puppe 序列、
悬挂、长正合列一路画下来，再对照链复形的公式
$d_{C_f}=\begin{pmatrix}d_A[1]&0\\ f&d_B\end{pmatrix}$。最后用柱
$\cyl(A\udot)$ 验证“同伦映射在 $D(\A)$ 中相等”，这是把拓扑直觉翻译成代数的
关键一步。

\paragraph{§\ref{sec:df} 导出函子。}把“适应类”三条当成消解存在的公理来读；
然后记住 $\R F$ 的复合定义与两个好处（复形是不变量；导出函子就是原函子）。
例 \ref{eg:derham} 是全篇的收束：de Rham 定理 = 常值层 $\mathbb R$ 的导出函子。
注意原文此处的 Dolbeault 命名与 Hodge 指标问题（校勘 4）。

\subsection*{D.3 五个关键论证拆解}

\paragraph{论证一：用锥证明 Whitehead 定理。}
设 $f\colon X\to Y$ 诱导同调同构。长正合列
$H_i(X)\to H_i(Y)\to H_i(C_f)\to H_{i-1}(X)$ 中，两端映射都是同构，
故中间项 $H_i(C_f)=0$ 对一切 $i$ 成立。由 Hurewicz 定理（单连通，
逐次上升）得 $\pi_*(C_f)=0$；再由 Whitehead 定理（$\pi_*$ 同构 $\Rightarrow$
同伦等价）得 $C_f\simeq$ 点。最后由 $C_f$ 的构造读出 $f$ 是同伦等价。
\emph{要点}：锥把“$f$ 的障碍”集中到一个空间 $C_f$ 上，于是“$f$ 是同构”
等价于“$C_f$ 可缩”。

\paragraph{论证二：链复形锥的两种上同调。}
对 $f\colon A\to B$（集中在 $0$ 次），$C_f=\{A\xrightarrow{f}B\}$：
$h^0=\ker f$、$h^1=\operatorname{coker}f$。这说明锥\emph{同时}记录核与余核，
只是分居不同次数。对一般的链映射，锥就是二项复形的全复形；证明
$C_i\simeq A[1]$ 只需写出显式的链映射并检查它在同调上是同构。

\paragraph{论证三：局部化的分数演算。}
为什么每个态射都能写成 $gs^{-1}$？因为拟同构对拉回封闭：给定
$X\xleftarrow{s}Z\xrightarrow{g}Y$ 与 $X\xleftarrow{s'}Z'$，可以拉回得到
$W$ 同时到 $Z$、$Z'$，再用消去条件把两个屋顶化成一个。这就是 Ore 条件；
它保证“形式地添逆”不会产生无穷长的词。参见 Gelfand--Manin III 2.8。

\paragraph{论证四：$D(\A)$ 为什么没有核与余核。}
若 $D(\A)$ 有核，则对 $0\to A\to B$ 应得到 $A$；但把 $A$ 换成拟同构的
$A'$ 后，核应当同构，而“逐项取核”并不保持拟同构。三角范畴用“补齐成特异三角”
替代核：$A\to B$ 的“余核”是 $C$，而 $C$ 只定义到同构，不典范——这正是
§\ref{sec:df} 末尾警告的来源。

\paragraph{论证五：de Rham 定理 = 导出函子。}
常值层 $\mathbb R$ 有消解 $A^\bullet(\mathbb R)$（Poincaré 引理），
而 $A^i(\mathbb R)$ 是细层，对 $\Gamma$ 无环。于是
$\R\Gamma(\mathbb R)\cong\Gamma(A^\bullet(\mathbb R))=\Omega^\bullet(X)$，
取上同调得 $H^*_{\mathrm{dR}}(X)\cong H^*(X;\mathbb R)$。
\emph{要点}：无环性保证“消解可以直接作用 $\Gamma$”，这就是“适应类”第二条
公理的用武之地。

\subsection*{D.4 与后续发展的关系}

本文停在“锥不能函子化选取”这一警告上。此后的发展大致三条线：
(1) \textbf{dg 范畴与 $A_\infty$ 范畴}——把 Hom 复形化，锥可以函子化
（Bondal--Kapranov、Kontsevich）；
(2) \textbf{稳定 $\infty$-范畴}——把“锥只定义到同构”提升为同伦 coherent 的
结构（Lurie）；
(3) \textbf{导出代数几何}——把导出范畴当作基本对象，反过来定义空间
（Toën--Vezzosi、Lurie）。
读本文时不必展开这些，但要知道：作者最后那句“还应当有某种更精细的概念”，
在二十年后已经成为主流。

\subsection*{D.5 五道习题（附提示）}

\begin{enumerate}[leftmargin=1.8em,itemsep=4pt]
\item 证明：若 $f\colon A\to B$ 是链映射，则 $C_f$ 的平移
  $C_f[1]$ 与 $C_{f[1]}$ 同构。\\
  \emph{提示}：写出 $d[1]=-d$ 的符号规则，逐项验证。
\item 验证 §\ref{sec:cone} 中的柱公式
  $d=\begin{pmatrix}d_A&0&0\\-\id&d_A[1]&\id\\0&0&d_A\end{pmatrix}$
  满足 $d^2=0$。\\
  \emph{提示}：注意 $d_A[1]=-d_A$，三处符号正好相消。
\item 设 $A=\C[\epsilon]/(\epsilon^2)$，证明
  $\R\mathcal Hom_A(k,A)\cong k[1]$，并由此说明
  $k\otimes^L_A\R\mathcal Hom_A(k,A)\not\cong\R\mathcal Hom_A(k,k)$。\\
  \emph{提示}：用消解 $A\xrightarrow{\epsilon}A$ 计算。
\item 对除子 $D\subset X$，证明
  $\Tor_1(\OO_D,\OO_D)\cong L^{-1}|_D$，并解释为什么
  $c_1(\OO_D)-c_1(\Tor_1)=D\cdot D$。\\
  \emph{提示}：用 \eqref{eq:div} 作平坦消解；用 $c_1(L^{-1}|_D)=-D\cdot D$。
\item 举例说明：两个复形同调相同，但在 $D(\A)$ 中不同构。\\
  \emph{提示}：习题 3 的 $k$ 与 $k[1]$ 已给出一个；再试
  $A=k[x]/(x^2)$、复形 $A\xrightarrow{x}A$ 与 $k\oplus k[1]$。
\end{enumerate}
